\section{Belief-Augmented Monte Carlo Search}
\label{sec:bams_search}

\subsection{Asymmetric 2-Ply Nested Search Formulation}

To overcome the computational cost of deep tree search while eliminating the strategy fusion flaw of PIMC, \algo\ implements \textbf{Belief-Augmented Monte Carlo Search (\bams)}. Instead of simulating deep unguided trajectories, \bams\ evaluates candidate plays across a structured, asymmetric 2-ply nested lookahead:

\begin{enumerate}
    \item \textbf{Root Candidate Beam ($K = 15$):} The native C GADDAG generator generates all legal plays for current rack $\rack_1$ on board $\board$, pruning down to the top-$K$ candidate moves ranked by initial static equity $\score(c_i) + \leaveval(\rack_1 \setminus c_i)$.
    \item \textbf{Opponent Counterplay Modeling ($N = 24$):} For each candidate play $c_i$, the hypothetical child board $\board_i = \board \oplus c_i$ is constructed. \bams\ draws $M$ opponent rack samples $\rack_{\opp}^{(m)} \sim b_t(\rack_{\opp})$ from the active Bayesian belief distribution. For each sample, the generator produces top-$N$ opponent replies $o_j \in \cA(\board_i, \rack_{\opp}^{(m)})$. The opponent's optimal counterplay $\hat{o}_i^{(m)}$ is identified:
    \begin{equation}
        \hat{o}_i^{(m)} = \arg\max_{o_j} \Big[ \score(o_j) + \leaveval(\rack_{\opp}^{(m)} \setminus o_j) \Big]
    \end{equation}
    \item \textbf{Player Reply Evaluation:} Conditioning on opponent counterplay $\hat{o}_i^{(m)}$ and our simulated bag draw $d \sim \mathrm{Hypergeometric}(\bag \setminus \rack_{\opp}^{(m)}, |c_i|)$, the child state transitions to:
    \begin{equation}
        \board_i' = \board_i \oplus \hat{o}_i^{(m)}, \quad \rack_1' = (\rack_1 \setminus c_i) \cup d
    \end{equation}
    Our best follow-up response $\hat{r}_i^{(m)} \in \cA(\board_i', \rack_1')$ is evaluated using the neural value network $\hat{V}$.
    \item \textbf{Expected 2-Ply Return:} The net equity of candidate $c_i$ is computed as:
    \begin{equation}
        \label{eq:bams_equity}
        Q_{\bams}(c_i) = \score(c_i) + \frac{1}{M} \sum_{m=1}^M \Big[ -\score(\hat{o}_i^{(m)}) + \score(\hat{r}_i^{(m)}) + \hat{V}(\board_i', \rack_1'') \Big]
    \end{equation}
\end{enumerate}

\begin{figure}[t]
\begin{lstlisting}[language=Python,caption={High-Level Algorithmic Structure of BAMS Lookahead}]
def bams_evaluate_play(board, rack, cand_play, belief_dist, M=16, N=24):
    board_after = apply_move(board, cand_play)
    tiles_played = cand_play.tiles_played
    rollout_scores = []
    
    for m in range(M):
        # Sample opponent rack from Bayesian belief posterior
        opp_rack = belief_dist.sample_rack()
        # Find best opponent counterplay on board_after
        opp_cand = generator.generate_plays(board_after, opp_rack, top_n=N)
        best_opp = max(opp_cand, key=lambda p: p.score + p.leave_val)
        
        # Sample player replenishment draw
        draw = sample_hypergeometric(bag_unseen, len(tiles_played))
        my_next_rack = (rack - tiles_played) + draw
        
        # Board after opponent counterplay
        board_next = apply_move(board_after, best_opp)
        my_replies = generator.generate_plays(board_next, my_next_rack, top_n=5)
        best_reply = max(my_replies, key=lambda p: p.score + value_net(p))
        
        # Net 2-ply outcome
        rollout_scores.append(cand_play.score - best_opp.score + best_reply.score)
        
    return mean(rollout_scores)
\end{lstlisting}
\end{figure}

\subsection{Bayesian Negative-Evidence Opponent Belief Tracking}

Classical engines assume unplayed opponent tiles are uniformly distributed across all unseen tiles $U = \bag \cup \rack_2$:
\begin{equation}
    \rack_2 \sim \mathrm{Hypergeometric}(U, 7)
\end{equation}
However, in expert human and AI play, unplayed letters leave definitive behavioral footprints. When an open bingo lane yielding $\ge 60$ points exists on the board, and the opponent plays a modest 14-point play elsewhere, the hypothesis that the opponent holds power bingo tiles ($?$, S) becomes exponentially improbable.

\algo\ models opponent rack distributions via a Bayesian tracker:
\begin{equation}
    b_{t+1}(\rack_2) \propto b_t(\rack_2) \cdot \Prob(a_{\opp} \mid \board_t, \rack_2)
\end{equation}
We formulate the likelihood function $\Prob(a_{\opp} \mid \board_t, \rack_2)$ via negative evidence:
\begin{equation}
    \Prob(a_{\opp} \mid \board_t, \rack_2) = \begin{cases}
        1.0 & \text{if } a_{\opp} \text{ was top play}, \\
        \exp(-\Delta_{\mathrm{forf}} / T) & \text{otherwise}
    \end{cases}
\end{equation}
where $\Delta_{\mathrm{forf}}$ is the forfeited point spread and $T = 12.5$ is a calibrated behavioral temperature. This tracker eliminates ``ghost threats'': \algo\ recognizes when an opponent is stuck with unbalanced vowel racks or lacks blank tiles, permitting bold offensive lane creation without defensive paralysis.

\subsection{Volatility-Conditioned Margin Control}

Point differential $\Delta = S_1 - S_2$ fundamentally alters risk preferences in Scrabble. Holding a 60-point lead in the midgame dictates board closure and variance reduction; conversely, trailing by 60 points requires high-volatility play to create mutual bingo opportunities.

\algo\ modulates candidate scoring using a continuous volatility penalty:
\begin{equation}
    \label{eq:volatility}
    \mathrm{Score}_{\mathrm{adj}}(c_i) = Q_{\bams}(c_i) - \lambda(\Delta) \cdot \Var_{\omega}[\score(\hat{o}_i)]
\end{equation}
where the risk weight $\lambda(\Delta)$ ramps continuously across margins:
\begin{equation}
    \lambda(\Delta) = \mathrm{clamp}\left(0.15 + \frac{\Delta}{100.0}, -0.30, +0.60\right)
\end{equation}
When leading ($\Delta > 0$), $\lambda > 0$, penalizing plays with high opponent response variance. When trailing ($\Delta < 0$), $\lambda < 0$, rewarding plays that widen opponent response variance and break open locked board configurations.
